\section*{第 \ref{ch:g6:wholes} 章——整数}

\begin{solution}{exo:g6:wholes:1}
$3\,047$; \quad $20\,500$; \quad $104$. \quad
$60\,013$：“六万零十三”。
\end{solution}

\begin{solution}{exo:g6:wholes:2}
百位数字是 $2$，个位数字是 $9$；$4$ 是
万位上的数字。分解如下：
\[
47\,209 = 4 \times 10\,000 + 7 \times 1000 + 2 \times 100 + 0 \times 10
+ 9 .
\]
\end{solution}

\begin{solution}{exo:g6:wholes:3}
$507 < 570$（十位：$0 < 7$）；\quad
$1\,099 > 989$（四位数比三位数大）；\quad
$23\,456 < 23\,465$（十位：$5 < 6$）；\quad
$9\,999 < 10\,000$（四位数与五位数相比）。
\end{solution}

\begin{solution}{exo:g6:wholes:4}
按位数比较，$978$ 和 $1\,001$ 排在
首位为 $2$ 的三个四位数之前；在这三个数中，比较十位和
个位：$2\,034 < 2\,043 < 2\,304$。最终顺序为：
\[
978 < 1\,001 < 2\,034 < 2\,043 < 2\,304 .
\]
\end{solution}

\begin{solution}{exo:g6:wholes:5}
$348 + 576 = 924$；\quad $805 - 268 = 537$（验算：
$537 + 268 = 805$); \quad $307 \times 26 = 307 \times 20 + 307 \times 6
= 6\,140 + 1\,842 = 7\,982$.
\end{solution}

\begin{solution}{exo:g6:wholes:6}
$38 + 45 + 12 + 5 = (38 + 12) + (45 + 5) = 50 + 50 = 100$.

$2 \times 83 \times 50 = (2 \times 50) \times 83 = 100 \times 83 =
8\,300$.

$4 \times 78 \times 25 = (4 \times 25) \times 78 = 100 \times 78 =
7\,800$.
\end{solution}

\begin{solution}{exo:g6:wholes:7}
$95 = 7 \times 13 + 4$（商为 $13$，\omterm{def:g3:division:remainder}{余数}为 $4$）。

$230 = 6 \times 38 + 2$（商为 $38$，\omterm{def:g3:division:remainder}{余数}为 $2$）。

$504 = 8 \times 63 + 0$（商为 $63$，\omterm{def:g3:division:remainder}{余数}为 $0$：$504$ 能
被 $8$ \omterm{def:g6:wholes:divisible}{整除}）。
\end{solution}

\begin{solution}{exo:g6:wholes:8}
等式成立，但\omterm{def:g3:division:remainder}{余数} $13$ \emph{不}小于
除数 $8$，因此它不满足带余除法的要求。再从\omterm{def:g3:division:remainder}{余数}中取出一个
$8$，得 $173 = 8 \times 21 + 5$，其中 $0 \leq 5 < 8$：
商为 $21$，\omterm{def:g3:division:remainder}{余数}为 $5$。
\end{solution}

\begin{solution}{exo:g6:wholes:9}
\emph{1.} $2\,000 = 12 \times 166 + 8$：农场装满了 $166$ 盒
，还剩 $8$ 个鸡蛋。

\emph{2.} $150 = 12 \times 12 + 6$：十二盒只能装 $144$
个鸡蛋，不够。面包店需要买 $13$ 盒。
\end{solution}

\begin{solution}{exo:g6:wholes:10}
$100 = 7 \times 14 + 2$：一百天是 $14$ 个整周，再加 $2$
天。十四周之后又是星期二；星期二之后的两天是
\emph{星期四}。
\end{solution}

\begin{solution}{exo:g6:wholes:11}
用 $9$ 去除时，最大的可能\omterm{def:g3:division:remainder}{余数}是 $8$（\omterm{def:g3:division:remainder}{余数}
必须小于除数）。因此被除数为
\[
9 \times 56 + 8 = 504 + 8 = 512 .
\]
\end{solution}

\begin{solution}{exo:g6:wholes:12}
设百位、十位、个位上的数字分别为 $h$、$0$、$u$。已知
$h = 2 \times u$，且 $h + 0 + u = 9$，于是 $2u + u = 9$，得
$3u = 9$，$u = 3$，$h = 6$。这个数是 $603$。验算：$6$ 是
$3$ 的两倍，十位数字是 $0$，且 $6 + 0 + 3 = 9$。
\end{solution}

\begin{solution}{pb:g6:wholes:1}
\textbf{1.} 例如 $(1 + 4) + (2 + 3) + 5 = 5 + 5 + 5 =
15$.

\textbf{2.} 每一对的和都是 $11$：$1 + 10$、$2 + 9$、$3 + 8$、
$4 + 7$、$5 + 6$。共 $5$ 对（十个数，每对两个数），
所以总和为 $5 \times 11 = 55$。

\textbf{3.} 每对之和为 $21$（$1 + 20$、$2 + 19, \dots$），并且
共有 $10$ 对：$10 \times 21 = 210$。

\textbf{4.} 在每一对中，第一个数增加 $1$ 时
第二个数就减少 $1$，因此总和始终不变：
$1 + 100 = 2 + 99 = 3 + 98 = \dots = 101$。这一百个数
恰好组成 $50$ 对（每对两个数），最内层的一对是
$50 + 51 = 101$.

\textbf{5.} $50$ 对，每对之和为 $101$：
\[
1 + 2 + \dots + 100 = 50 \times 101 = 5\,050 .
\]

\textbf{6.} 九个数中，中间的数 $5$ 没有
伙伴。它周围有 $1 + 9 = 2 + 8 = 3 + 7 = 4 + 6 = 10$，共
四对。总和为 $4 \times 10 + 5 = 45$。

\textbf{7.} 每一列之和为 $1 + 9 = 2 + 8 = \dots =
9 + 1 = 10$，共有 $9$ 列，两行合起来
之和为 $9 \times 10 = 90$。但这两行是
\emph{同一个}和写了两次，所以其中一行之和恰为\omterm{def:g3:division:half}{一半}：
$1 + 2 + \dots + 9 = 90 \div 2 = 45$。没有任何一个数会
落单——它的伙伴正好在它下面。

\textbf{8.} 对于 $1$ 到 $50$：$50$ 列，每列之和为 $51$，两倍的和
为 $50 \times 51 = 2\,550$，因此和为 $2\,550 \div 2 = 1\,275$。
对于 $1$ 到 $200$：$200$ 列，每列之和为 $201$，两倍的和
为 $40\,200$，因此和为 $20\,100$。

\textbf{9.} 每个\omterm{def:g3:numbers:evenodd}{偶数}都是高斯和中某个数的两倍：
$2 = 2 \times 1$，$4 = 2 \times 2$，\dots，
$200 = 2 \times 100$。因此整个和是高斯和的两倍：
\[
2 + 4 + \dots + 200 = 2 \times 5\,050 = 10\,100 .
\]

\textbf{10.} $101$ 到 $200$ 这些数，就是
$1$ 到 $200$ 的数去掉 $1$ 到 $100$ 的数：
\[
101 + 102 + \dots + 200 = 20\,100 - 5\,050 = 15\,050 .
\]

\textbf{11.} 把朋友们排成一列。第一个人握了 $7$ 次手，
分别与其他每人各握一次。第二个人已经和第一个人打过招呼，所以只需
\emph{新}握 $6$ 次；第三个人握 $5$ 次；依此类推，第七个人
握 $1$ 次，第八个人则没有新的握手。
每次握手恰好只计算一次：
\[
7 + 6 + 5 + 4 + 3 + 2 + 1 = 28 \text{ 次握手}
\]
（围绕中间配对：$(7+1) + (6+2) + (5+3) + 4 =
8 + 8 + 8 + 4 = 28$）。

\textbf{12.} $1 + 2 + \dots + 12$：十二列，每列之和为 $13$，
两倍的和为 $12 \times 13 = 156$，因此
$156 \div 2 = 78$ 罐。

\textbf{13.} 由 $1 + 2 + 3 + 4$ 个方块排成的阶梯，再加上
同一个倒置的阶梯，恰好拼成一个 $4$ 行 $5$
列的长方形：长方形的每一行对应两行配对法中的一列，
（上升阶梯的一级与下降阶梯的一级配成一对，
即 $1 + 4$、$2 + 3$、$3 + 2$、$4 + 1$——
总是 $5$）。因此两倍的和为 $4 \times 5 = 20$ 个方块，而
$1 + 2 + 3 + 4 = 10$：此图\emph{正是}第 7 题。

\textbf{14.} $1 + 2 + \dots + 12 = 78$，表示十二小时内钟声
敲响的次数（见第 12 题），因此一整天共敲响 $78 \times 2 = 156$ 次。

\textbf{15.} 除以 $60$：$5\,050 = 60 \times 84 + 10$，因此
$5\,050$ 秒是 $84$ 分钟又 $10$ 秒。将
分钟数除以 $60$：$84 = 60 \times 1 + 24$，因此 $84$ 分钟是
$1$ 小时 $24$ 分钟。合计 $1$ h $24$ min $10$ s——略
\emph{少于}一个半小时（$1$ h $30$ min）。
\end{solution}
